\chapter{Projection Lemmas}

Let $\mathcal{A}$ be a unital $C^{*}$--algebra in this section. In this section we collect relevant results about (selfadjoint) projections in $\mathcal{A}$.
Recall that an element $p$ in a $C^{*}$--algebra is a projection if $p^2=p^{*}=p$. By the Spectral Mapping Theorem (or using the full Gelfand Duality),
a selfadjoint element $a\in A$ is a projection if and only if the spectrum of $a$ is contained in $\{0,1\}$.

\begin{lemma}
  \label{lem:proj_le_one}
  \lean{IsStarProjection.le_one}
  \mathlibok
  More generally, every projection $p$ in a star-ordered ring satisfies $p\le 1$.
\end{lemma}

\begin{proof}
  \leanok
  The complement $1-p$ is again a projection and is therefore nonnegative. This is exactly
  $p\le 1$.
\end{proof}

\begin{lemma}
  \label{lem:proj_sub_pos_iff_comm_eq_self}
  \lean{IsStarProjection.le_tfae}
  \leanok
  \uses{lem:star_conj_pos, lem:proj_le_one}
  Let $\mathcal{A}$ be a $C^{*}$--algebra and $p,q \in \mathcal{A}$ be projections. Then $p-q\ge 0$ iff $qp = pq = q$.
\end{lemma}

\begin{proof}
  \leanok
  By Lemma \ref{lem:star_conj_pos}, if $p-q \ge 0$ then $q(p-q)q=qpq-q^3=qpq-q \ge 0$. By Lemma \ref{lem:proj_le_one}, we have $p\le 1$ and
  employing Lemma \ref{lem:star_conj_pos} we obtain $qpq\le q1q=q$. Since $qpq\ge q$ and $qpq\le q$ we have $qpq=q$, which implies that $q(p-q)q=0$.
  By the $C^{*}$-property of the norm, we have $\|(p-q)^{1/2}q\|^2=\|q(p-q)q\|=0$ and so $(p-q)^{1/2}q=0$ and therefore $(p-q)q=(p-q)^{1/2}(p-q)^{1/2}q=0$.
  It follows that $pq=q$, and taking adjoints, $qp=q$.
  Conversely, if $qp=pq=q$, one easily checks that $p-q$ is a projection and so its spectrum is contained in $\{0,1\}$ and it is positive.
\end{proof}

\begin{corollary}
  \label{cor:proj_sub_of_subproj}
  \lean{IsStarProjection.le_tfae}
  \leanok
  \uses{lem:proj_sub_pos_iff_comm_eq_self}
  For all $p,q \in \mathcal{A}$ projections such that $q\le p$, $p-q$ is a projection.
\end{corollary}
\begin{proof}
  \leanok
\end{proof}

The next results are specific to $W^{*}$--algebras.

\begin{proposition}
  \label{prop:left_ideals_in_wstar_proj_Sak_1_10_1} (Sakai 1.10.1, left-ideal case)
  \lean{Ideal.existsUnique_isStarProjection_eq_span_of_isClosed_ultraweak,CStarAlgebra.isUnital_of_isClosed_submodule}
  \leanok
  \uses{lem:star_left_mul_right_mul_sig_cts_Sak_1_7_8,lem:wstar_unital}
  Let $\mathcal{L}$ be a $\sigma(M,M_{*})$--closed left ideal of a $W^{*}$--algebra $M$.
  Then there exists a unique projection $p$ in $M$ so that $\mathcal{L}=Mp$.
  In Lean, left ideals use Mathlib's native type \texttt{Ideal M}, and $Mp$ is represented by
  \texttt{Ideal.span \{p\}}.
\end{proposition}

\begin{proof}
  \leanok
  Suppose $\mathcal{L}$ is a $\sigma$--closed left ideal. By Lemma \ref{lem:star_left_mul_right_mul_sig_cts_Sak_1_7_8}, $\mathcal{L^{*}}$ is
  also $\sigma$--closed, so $\mathcal{N}=\mathcal{L}\cap \mathcal{L}^{*}$ is a $\sigma$--closed
  nonunital $*$--subalgebra of $M$. Its explicit closed-submodule predual, together with Lemma
  \ref{lem:wstar_unital}, supplies a multiplicative identity $p$ for $\mathcal{N}$; this identity is
  a selfadjoint projection in $M$. Certainly $Mp\subseteq\mathcal{L}$. Conversely, if
  $x\in\mathcal{L}$, then $x^{*}x\in\mathcal{N}$ and hence $(x^{*}x)p=x^{*}x$. The $C^{*}$--identity
  applied to $x-xp$ gives $xp=x$, so $\mathcal{L}=Mp$.

  If projections $p$ and $q$ generate the same left ideal, the two ideal inclusions give
  $pq=p$ and $qp=q$. Taking adjoints of the first equality also gives $qp=p$, and therefore
  $p=q$.
\end{proof}

\begin{corollary}
  \label{cor:right_ideals_in_wstar_proj_Sak_1_10_1} (Sakai 1.10.1, right-ideal case)
  \lean{Ideal.existsUnique_isStarProjection_eq_span_op_of_isClosed_ultraweak}
  \leanok
  \uses{prop:left_ideals_in_wstar_proj_Sak_1_10_1}
  Let $\mathcal{R}$ be a $\sigma(M,M_{*})$--closed right ideal of a $W^{*}$--algebra $M$.
  Then there exists a unique projection $q$ in $M$ so that $\mathcal{R}=qM$.
  Formally, a right ideal of $M$ is represented canonically as
  \texttt{Ideal (MulOpposite M)},
  and $qM$ is \texttt{Ideal.span \{MulOpposite.op q\}}.
\end{corollary}

\begin{proof}
  \leanok
  Transport the specified predual of $M$ to $M^{\mathrm{op}}$ along the linear isometry
  $\operatorname{unop}:M^{\mathrm{op}}\to M$. A right ideal of $M$ is then precisely a left ideal
  of $M^{\mathrm{op}}$, with the same $\sigma$--closedness condition. Apply Proposition
  \ref{prop:left_ideals_in_wstar_proj_Sak_1_10_1} in the opposite algebra and transport its unique
  projection back with $\operatorname{unop}$. The maps $\operatorname{op}$ and
  $\operatorname{unop}$ preserve star projections, and the principal left ideal generated by
  $\operatorname{op}(q)$ in $M^{\mathrm{op}}$ represents $qM$.
\end{proof}


\begin{proposition}
  \label{prop:proj_compl_lat_wstar_Sak_1_10_2} (Sakai 1.10.2)
  \lean{IsStarProjection.completeLatticeOfPredual,IsStarProjection.instCompleteLatticeSubtypeOfWStarAlgebra}
  \leanok
  \uses{prop:left_ideals_in_wstar_proj_Sak_1_10_1,lem:proj_sub_pos_iff_comm_eq_self,lem:star_left_mul_right_mul_sig_cts_Sak_1_7_8}
  Let $M$ be a $W^{*}$--algebra. Its set $M^{p}$ of projections is a complete lattice with respect
  to $\le$. In Lean, this is the subtype \texttt{\{p : M // IsStarProjection p\}}.
\end{proposition}

\begin{proof}
  \leanok
  Send a projection $p$ to its principal left ideal $Mp$. This ideal is ultraweakly closed, and
  Proposition \ref{prop:left_ideals_in_wstar_proj_Sak_1_10_1} gives a unique inverse on the
  ultraweakly closed left ideals. Lemma \ref{lem:proj_sub_pos_iff_comm_eq_self} identifies the two
  orders, so this is an order isomorphism.

  Arbitrary intersections of ultraweakly closed left ideals are again ultraweakly closed. Hence the
  infimum of a family of projections is the unique projection corresponding to the intersection of
  its principal left ideals. This supplies all infima; the standard order-theoretic construction
  then supplies the complete lattice. The explicit construction is parameterized by a specified
  predual. For the global $W^{*}$--algebra instance Lean uses the canonically selected predual; the
  resulting lattice operations are independent of that selection because infima and suprema are
  uniquely determined by the underlying projection order.
\end{proof}
